\section{Configuración Experimental y Desarrollo del Prototipo}
\label{sec:configuracion_experimental_prototipo}

Esta sección expone los aspectos metodológicos e instrumentales relativos a la implementación del prototipo experimental y la infraestructura de pruebas. Se describe la plataforma de computación, el uso de las APIs de Google Gemini, la base de datos vectorial Qdrant y las salvaguardas arquitectónicas implementadas para operar bajo restricciones de cuota estrictas. A continuación, se detallan las decisiones de ingeniería software aplicadas en los módulos de ingesta, vectorización, reordenación bimodal e interfaz gráfica. Finalmente, se presenta el diseño del arnés de auditoría y trazabilidad (\textit{RunAuditor}) que asegura la total reproducibilidad científica de los resultados empíricos.

\subsection{Entorno Experimental, Infraestructura y Gestión de Cuotas}
\label{subsec:entorno_experimental_infraestructura}

Esta subsección describe la plataforma de cómputo, las librerías del entorno de ejecución y las políticas de control de cuotas y memoria del prototipo.

\subsubsection{Plataforma de Ejecución y Pila Tecnológica}
Los experimentos y el desarrollo del prototipo se ejecutaron sobre una estación de trabajo con procesador Apple Silicon M3 (arquitectura ARM64), $16\text{ GB}$ de memoria unificada y sistema operativo macOS, utilizando el entorno de ejecución Python 3.13 (\texttt{3.13.12}). La pila de software científico y librerías especializadas comprende:
\begin{itemize}
    \item \textbf{Procesamiento Audiovisual:} \texttt{OpenCV} (\texttt{opencv-python}) para la decodificación de metraje, cálculo matricial del índice SSIM, muestreo por entropía visual y compresión JPEG; \texttt{ffmpeg} para el remuestreo de audio; y \texttt{openai-whisper} (modelo \texttt{base}) para el reconocimiento automático del habla (ASR) ejecutado localmente en CPU.
    \item \textbf{Modelos Fundacionales y Motores de Vectorización (Cloud y Local):} SDK oficial \texttt{google-genai} para el acceso remoto a \textit{Gemini Embedding 2} (\texttt{gemini-embedding-2-preview}) y el árbitro bimodal \textit{Gemini 3.5 Flash Lite} (\texttt{gemini-3.5-flash-lite}); junto con \texttt{PyTorch} (\texttt{2.14.1}, acelerador Metal Performance Shaders \texttt{mps}), \texttt{transformers} (\texttt{5.22.2}), \texttt{sentence-transformers[image]} (\texttt{6.3.0}) y \texttt{torchvision} (\texttt{0.29.2}) para inferencia local de \textit{Qwen3-VL-Embedding-2B} y \textit{EmbeddingGemma 2} en precisión \texttt{bfloat16} sobre \texttt{mps} (y \texttt{float32} en CPU).
    \item \textbf{Indexación Vectorial:} \texttt{qdrant-client} en modo local persistente en disco (\texttt{./data/vector\_db}), evitando la sobrecarga de latencia de red y garantizando aislamiento experimental entre configuraciones mediante colecciones independientes.
    \item \textbf{Cálculo Numérico, Estadístico y Evaluación:} \texttt{NumPy} para el truncamiento dimensional y normalización $L_2$ bajo el estándar MRL, \texttt{SciPy} (\texttt{1.18.1}) para las pruebas de significación pareada (\textit{McNemar} y \textit{Wilcoxon}), y \texttt{pytest} como motor de ejecución para las suites de prueba unitarias y de extremo a extremo.
\end{itemize}

\subsubsection{Gestión de Cuotas, Concurrencia y Presupuesto de Memoria Unificada}
La capa gratuita de la API de Google Gemini impone restricciones severas que condicionan el diseño de cualquier sistema experimental en tiempo real: un máximo de $100$ peticiones por minuto (RPM), $30.000$ tokens por minuto (TPM) y $1.000$ peticiones diarias (RPD) \cite{gemini_emb2_docs}. En paralelo, la ejecución local de modelos multimodales de hasta $2\text{B}$ parámetros dentro de $16\text{ GB}$ de RAM unificada requiere evitar la fragmentación de caché en GPU (\texttt{mps}) y las colisiones concurrentes de descriptores de comandos.

Para operar de manera robusta y continua sin interrupciones por denegación de servicio (código HTTP 429 / \texttt{RESOURCE\_EXHAUSTED}) ni desbordamiento de memoria, se implementaron cuatro salvaguardas arquitectónicas:
\begin{enumerate}
    \item \textbf{Reintentos Exponenciales con Dispersión (\textit{Exponential Backoff with Jitter}):} Toda llamada remota a la API se canaliza a través de la función utilitaria \texttt{call\_with\_retry}. Ante errores transitorios de cuota o sobrecarga de red, el cliente aplica una pausa adaptativa con retardo base de $4,0\text{ s}$ para vectorización y $1,5\text{ s}$ para reordenación, incrementándose exponencialmente con un factor de aleatoriedad hasta un máximo de 4 reintentos antes de elevar una excepción controlada (\texttt{VectorizationError}).
    \item \textbf{Cortafuegos de Cuota por Vídeo (\texttt{QuotaSafetyError}):} Antes de vectorizar, el orquestador verifica que el total de fragmentos no supere $\text{MAX\_CHUNKS\_PER\_VIDEO} = 800$. Si un vídeo supera dicho umbral, aborta preventivamente antes de emitir llamadas a la API, evitando agotar la cuota diaria.
    \item \textbf{Semáforo Adaptativo y Cerrojo por Modelo:} En \texttt{process\_\allowbreak video\_\allowbreak pipeline}, las llamadas remotas a Gemini emplean un máximo de $3$ peticiones concurrentes (\texttt{max\_\allowbreak concurrent\_\allowbreak requests = 3}), mientras que los motores locales fuerzan concurrencia estricta de $1$ (\texttt{Semaphore(1)}) y serializan las pasadas \texttt{model.encode} mediante un \texttt{threading.Lock} compartido por \texttt{model\_id} en \texttt{local\_\allowbreak vectorizer.py}.
    \item \textbf{Descarga Secuencial de Modelos (\textit{Single-Resident Lifecycle}):} Antes de cargar un vectorizador local tras la transcripción ASR, el sistema invoca \texttt{unload\_whisper\_models()}; de manera simétrica, la factoría \texttt{factory.py} mantiene un único modelo local residente y libera el anterior mediante \texttt{unload\_active\_vectorizer()} (\texttt{gc.collect()} seguido de \texttt{torch.mps.empty\_cache()}).
\end{enumerate}

\subsubsection{Configuración y Aislamiento de Colecciones en Qdrant}
La base vectorial local almacena los fragmentos indexados en colecciones aisladas generadas por \texttt{collection\_name\_for(backend, dim)}. Para preservar el índice histórico de referencia sin consumir cuota diaria adicional, la configuración predeterminada (\texttt{gemini} en $d=3072$) conserva la colección \texttt{video\_chunks}, mientras que cualquier otra combinación de motor y dimensionalidad se aísla bajo el esquema determinista \texttt{video\_chunks\_\_<backend>\_\_<dim>} (cubriendo \texttt{gemini} en $1536/768$, \texttt{qwen3-vl-2b} en $2048/1536/768$ y \texttt{gemma2} en $768/512/256$). Cada colección se configura con métrica de distancia coseno (\texttt{Distance.COSINE}) y almacena un \textit{payload} documental uniforme (nombre de vídeo, \texttt{chunk\_index}, intervalo $[t_{\text{start}}, t_{\text{end}}]$ y transcripción), habilitando la conmutación instantánea de motores en la interfaz sin reindexar.

\newpage

\subsection{Implementación Modular del Prototipo Experimental}
\label{subsec:implementacion_modular_prototipo}

El prototipo se diseñó siguiendo el principio de inversión de dependencias y bajo un desacoplamiento estricto de componentes, como evidencia el análisis de grafos del repositorio (27 comunidades funcionales y 0 ciclos de dependencia). A continuación se detallan las decisiones de implementación adoptadas en los módulos nucleares:

\subsubsection{Módulo de Ingesta, Segmentación y Transcripción}
El paquete \texttt{src/processing/} implementa la ingesta de medios a través de dos componentes cooperativos:
\begin{itemize}
    \item \textbf{\texttt{video\_chunker.py}:} La función \texttt{process\_video\_single\_pass} recorre los fotogramas del metraje decodificados por OpenCV. En lugar de almacenar fotogramas completos en memoria, calcula el índice SSIM sobre miniaturas en escala de grises de 160 píxeles de ancho (\texttt{resize\_width = 160}), acelerando el cálculo en un factor de $8\times$. Cuando se detecta un corte de plano, extrae los tres fotogramas clave en su resolución de trabajo (ancho 512 px) y los almacena en una caché local temporal en formato JPEG para evitar re-lecturas de disco.
    \item \textbf{\texttt{audio\_transcriber.py}:} Ejecuta Whisper en un proceso aislado, capturando el flujo de segmentos temporales. La función \texttt{align\_transcript\_with\_chunks} aplica el algoritmo de solape temporal ($\ge 0,2\text{ s}$) para fusionar de manera forzada las locuciones con cada \texttt{VideoChunk}.
\end{itemize}

\subsubsection{Vectorizador Multimodal y Truncamiento MRL}
El módulo \texttt{src/embeddings/gemini\_vectorizer.py} define la clase \texttt{GeminiMultimodalVectorizer}. Este componente ensambla el objeto \texttt{types.Content} combinando las tres imágenes JPEG codificadas en binario con la parte textual de la transcripción fonética.

El truncamiento de representaciones Matryoshka se encapsula en la función matemática \texttt{truncate\_and\_normalize\_mrl}:
\begin{equation}
    \mathbf{v}^{(d)} = \text{\texttt{truncate\_and\_normalize\_mrl}}(\mathbf{v}, d)
\end{equation}
Esta función comprueba que la dimensión solicitada $d$ pertenezca a los límites admisibles ($1 \le d \le 3072$), extrae la sub-rebanada (\textit{slice}) de los primeros $d$ coeficientes, calcula su norma euclídea mediante \texttt{numpy.linalg.norm} y proyecta el vector resultante sobre la hiperesfera $\mathbb{S}^{d-1}$.

\subsubsection{Motor de Búsqueda Vectorial Qdrant}
La clase \texttt{QdrantSearchEngine} en \texttt{src/rag\_engine/vector\_db.py} gestiona la persistencia e indexación en la base de datos vectorial local. Para cada fragmento, genera un identificador determinista derivado del nombre del archivo y del índice secuencial del plano, evitando duplicados en re-indexaciones. La consulta de vecinos más cercanos se formula a través de la API \texttt{search} de Qdrant, extrayendo las cargas útiles asociadas a los vectores para reconstruir instancias de \texttt{SearchResult}.

\subsubsection{Árbitro Bimodal Multimodal y Parser Resiliente}
El componente \texttt{MultimodalReranker} en \texttt{src/rag\_engine/reranker.py} implementa la verificación bimodal de segunda etapa mediante llamadas asíncronas concurrentes. Para mitigar posibles truncamientos de tokens en respuestas JSON extensas emitidas por \textit{Gemini 3.5 Flash Lite}, se programó un decodificador sintáctico resiliente (\texttt{\_parse\_evaluations\_json}). Si el JSON estándar falla debido a un cierre incompleto de corchetes, el decodificador realiza un análisis iterativo mediante \texttt{json.JSONDecoder.raw\_decode}, recuperando todas las evaluaciones completas individuales antes del corte y evitando el fallo global de la consulta.

\subsubsection{Fachada de Orquestación e Interfaz Gráfica}
La clase orquestadora \texttt{VideoRAGService} (\texttt{src/rag\_engine/service.py}) actúa como fachada unificada del sistema, exponiendo dos métodos de alto nivel:
\begin{itemize}
    \item \texttt{index\_video(video\_path, run\_config, progress\_callback)}: ejecuta el pipeline completo de ingesta, transcripción, vectorización y almacenamiento en Qdrant.
    \item \texttt{search(query, search\_config, run\_config, auditor)}: orquesta la vectorización de la consulta multimodal, la recuperación en Qdrant, la deduplicación temporal y el reordenamiento bimodal.
\end{itemize}

Por último, el paquete \texttt{src/ui/app.py} implementa una interfaz gráfica reactiva con Streamlit. La aplicación permite componer consultas multimodales combinando cuadros de texto, carga de ficheros \texttt{.txt}, subida de imágenes y archivos de audio, visualizando los resultados mediante tarjetas interactivas que enlazan directamente el reproductor de vídeo con el intervalo temporal exacto $[t_{\text{start}}, t_{\text{end}}]$ recuperado.

\newpage

\subsection{Arnés de Trazabilidad, Observabilidad y Auditoría Experimental}
\label{subsec:arnes_trazabilidad_auditoria}

Uno de los pilares metodológicos exigidos en una investigación científica rigurosa es la auditabilidad completa de las pruebas empíricas. Los experimentos basados en modelos de lenguaje y APIs propietarias corren el riesgo de introducir sesgos de volatilidad temporal o pérdida de traza de las peticiones intermedias. Para blindar la validez científica del estudio, se implementó el sistema de auditoría integral \texttt{RunAuditor} (\texttt{src/ui/auditor.py} y \texttt{src/utils/audit.py}).

\subsubsection{Estructura del Registro Experimental}
Cada ejecución del sistema o batería de pruebas genera de manera determinista un directorio sellado temporalmente:
\begin{equation}
    \text{RutaAuditoría} = \text{\texttt{data/audit/}} \, \langle\text{timestamp}\rangle\_\langle\text{run\_id}\rangle /
\end{equation}

Dentro de esta ruta, el auditor persiste de forma atómica tres artefactos normalizados:
\begin{enumerate}
    \item \textbf{\texttt{run\_metadata.json}:} Registra la marca de tiempo de inicio y finalización en formato ISO-8601, el identificador único UUID del experimento, la versión del código y el desglose temporal granular: tiempo de decodificación de vídeo, tiempo de ejecución de Whisper, latencia agregada de llamadas a la API de embeddings, tiempo de consulta en el grafo HNSW de Qdrant y latencia de la etapa de verificación con el juez LLM.
    \item \textbf{\texttt{config.json}:} Contiene la serialización exhaustiva de todos los objetos de configuración inmutables utilizados (\texttt{ChunkingConfig}, \texttt{VectorizerConfig}, \texttt{SearchConfig} y \texttt{AudioConfig}). Esto garantiza que cualquier diferencia de rendimiento observada entre dos experimentos pueda atribuirse de manera inequívoca a una variación explícita en los hiperparámetros.
    \item \textbf{\texttt{search\_results.json}:} Almacena la traza acumulativa de cada consulta ejecutada durante la sesión. Para cada consulta, se archiva:
    \begin{itemize}
        \item Las modalidades presentes en la consulta de entrada (texto, imágenes o audios).
        \item La lista completa de candidatos devueltos por la Etapa 1 en Qdrant, incluyendo la puntuación de similitud coseno $S_C$, las marcas de tiempo y el texto transcrito de cada fragmento.
        \item Las decisiones emitidas por el árbitro bimodal en la Etapa 2: juicio booleano (\texttt{is\_relevant}), puntuación normalizada (\texttt{relevance\_score}) y la justificación textual del razonamiento emitido por el modelo (\texttt{reasoning}).
    \end{itemize}
\end{enumerate}

\subsubsection{Protocolo de Reproducibilidad Desacoplada}
Este diseño de trazabilidad introduce una ventaja científica determinante: desacopla por completo la fase de inferencia remota de la fase de análisis métrico.

Una vez ejecutada la batería de pruebas sobre un corpus, el archivo \texttt{search\_results.json} contiene la verdad experimental completa de la sesión. A partir de este registro, es posible ejecutar scripts de análisis cuantitativo automatizado para calcular métricas de recuperación (\textit{Recall@k}, \textit{MRR}, \textit{NDCG@k}, precisión de frontera temporal) y generar gráficos comparativos de forma instantánea, determinista y sin necesidad de realizar nuevas llamadas a la API ni incurrir en costes computacionales adicionales.

